\documentclass[main.tex]{subfiles}

\begin{document}

\section{Conclusion and References}
\begin{frame}
  \frametitle{Conclusion}
  \begin{enumerate}
    \item \textbf{AI workloads expose parallel linear algebra.}
      Neural-network layers repeatedly reduce to matrix and tensor operations.
    \item \textbf{GPUs already had the right foundation.}
      Graphics demanded parallel throughput and high memory bandwidth; unified shaders and CUDA made that hardware programmable.
    \item \textbf{Data movement determines real performance.}
      Coalescing, shared-memory tiling, and register reuse move GEMM toward the compute limit.
    \item \textbf{AI is now driving specialization.}
      Tensor Cores, lower precision, and new accelerators trade flexibility for throughput and efficiency.
  \end{enumerate}

  \begin{alertblock}{Main takeaway}
    Graphics and AI both reward massive parallelism. The GPU succeeded by evolving from a graphics pipeline into a programmable matrix engine.
  \end{alertblock}
\end{frame}

\begin{frame}
  \frametitle{References: GPU evolution}
  \small
  \begin{itemize}
    \item I. Buck et al. (2004), \href{https://graphics.stanford.edu/papers/brookgpu/brookgpu.pdf}{``Brook for GPUs: Stream Computing on Graphics Hardware.''}
    \item E. Lindholm et al. (2008), \href{https://www.cs.cmu.edu/afs/cs/academic/class/15869-f11/www/readings/lindholm08_tesla.pdf}{``NVIDIA Tesla: A Unified Graphics and Computing Architecture.''}
    \item NVIDIA (2009), \href{https://www.nvidia.com/content/PDF/fermi_white_papers/NVIDIA_Fermi_Compute_Architecture_Whitepaper.pdf}{\emph{Fermi Compute Architecture Whitepaper}.}
    \item NVIDIA (2017), \href{https://images.nvidia.com/content/volta-architecture/pdf/volta-architecture-whitepaper.pdf}{\emph{NVIDIA Tesla V100 GPU Architecture}.}
    \item NVIDIA (2020), \href{https://developer.nvidia.com/blog/nvidia-ampere-architecture-in-depth/}{``NVIDIA Ampere Architecture In-Depth.''}
    \item NVIDIA (2022), \href{https://developer.nvidia.com/blog/nvidia-hopper-architecture-in-depth/}{``NVIDIA Hopper Architecture In-Depth.''}
    \item NVIDIA (2024), \href{https://www.nvidia.com/en-us/data-center/technologies/blackwell-architecture/}{``NVIDIA Blackwell Architecture.''}
    \item NVIDIA, \href{https://www.nvidia.com/en-us/data-center/technologies/rubin/}{``NVIDIA Rubin Architecture.''}
  \end{itemize}
\end{frame}

\begin{frame}
  \frametitle{References: CUDA and GEMM}
  \small
  \begin{itemize}
    \item S. Boehm (2022), \href{https://siboehm.com/articles/22/CUDA-MMM}{``How to Optimize a CUDA Matmul Kernel for cuBLAS-like Performance.''}
    \item S. Williams, A. Waterman, and D. Patterson (2009), \href{https://doi.org/10.1145/1498765.1498785}{``Roofline: An Insightful Visual Performance Model for Multicore Architectures.''}
    \item NVIDIA, \href{https://docs.nvidia.com/cuda/cuda-c-programming-guide/}{\emph{CUDA C++ Programming Guide}.}
    \item NVIDIA, \href{https://www.nvidia.com/en-us/geforce/graphics-cards/30-series/rtx-3090-3090ti/}{``GeForce RTX 3090 Family Specifications.''}
    \item NVIDIA, \href{https://developer.nvidia.com/blog/cuda-refresher-cuda-programming-model/}{``CUDA Refresher: The CUDA Programming Model.''}
    \item NVIDIA, \href{https://docs.nvidia.com/cuda/cublas/}{\emph{cuBLAS Library Documentation}.}
    \item Modal, \href{https://modal.com/gpu-glossary/device-software/cuda-programming-model}{``CUDA Programming Model,'' \emph{GPU Glossary}.}
    \item F. Sanglard (2020), \href{https://fabiensanglard.net/cuda/}{``The CUDA Architecture.''}
  \end{itemize}
\end{frame}

\begin{frame}
  \frametitle{References: alternative accelerators}
  \footnotesize
  \begin{itemize}
    \item AMD, \href{https://www.amd.com/en/products/accelerators/instinct.html}{``AMD Instinct Accelerators.''}
    \item N. P. Jouppi et al. (2017), \href{https://doi.org/10.1145/3079856.3080246}{``In-Datacenter Performance Analysis of a Tensor Processing Unit.''}
    \item V. Sze et al. (2017), \href{https://doi.org/10.1109/JPROC.2017.2761740}{``Efficient Processing of Deep Neural Networks: A Tutorial and Survey.''}
    \item Cerebras, \href{https://www.cerebras.ai/chip}{``The Wafer-Scale Engine.''}
    \item Etched (2024), \href{https://www.etched.com/announcing-sohu}{``Introducing Sohu: The World's First Transformer ASIC.''}
  \end{itemize}

  \vfill
  \centering
  \large Questions?
\end{frame}

\end{document}
